% Options for packages loaded elsewhere
\PassOptionsToPackage{unicode}{hyperref}
\PassOptionsToPackage{hyphens}{url}
%
\documentclass[
  oneside, 11pt, a4paper]{memoir}
\usepackage{lmodern}
\usepackage{amssymb,amsmath}
\usepackage{ifxetex,ifluatex}
\ifnum 0\ifxetex 1\fi\ifluatex 1\fi=0 % if pdftex
  \usepackage[T1]{fontenc}
  \usepackage[utf8]{inputenc}
  \usepackage{textcomp} % provide euro and other symbols
\else % if luatex or xetex
  \usepackage{unicode-math}
  \defaultfontfeatures{Scale=MatchLowercase}
  \defaultfontfeatures[\rmfamily]{Ligatures=TeX,Scale=1}
  \setmainfont[]{Minion Pro}
  \setsansfont[]{Myriad Pro}
  \setmathfont[]{Minion Pro}
\fi
% Use upquote if available, for straight quotes in verbatim environments
\IfFileExists{upquote.sty}{\usepackage{upquote}}{}
\IfFileExists{microtype.sty}{% use microtype if available
  \usepackage[]{microtype}
  \UseMicrotypeSet[protrusion]{basicmath} % disable protrusion for tt fonts
}{}
\makeatletter
\@ifundefined{KOMAClassName}{% if non-KOMA class
  \IfFileExists{parskip.sty}{%
    \usepackage{parskip}
  }{% else
    \setlength{\parindent}{0pt}
    \setlength{\parskip}{6pt plus 2pt minus 1pt}}
}{% if KOMA class
  \KOMAoptions{parskip=half}}
\makeatother
\usepackage{xcolor}
\IfFileExists{xurl.sty}{\usepackage{xurl}}{} % add URL line breaks if available
\IfFileExists{bookmark.sty}{\usepackage{bookmark}}{\usepackage{hyperref}}
\hypersetup{
  pdfauthor={Januar Harianto},
  hidelinks,
  pdfcreator={LaTeX via pandoc}}
\urlstyle{same} % disable monospaced font for URLs
\usepackage{longtable,booktabs}
% Correct order of tables after \paragraph or \subparagraph
\usepackage{etoolbox}
\makeatletter
\patchcmd\longtable{\par}{\if@noskipsec\mbox{}\fi\par}{}{}
\makeatother
% Allow footnotes in longtable head/foot
\IfFileExists{footnotehyper.sty}{\usepackage{footnotehyper}}{\usepackage{footnote}}
\makesavenoteenv{longtable}
\usepackage{graphicx,grffile}
\makeatletter
\def\maxwidth{\ifdim\Gin@nat@width>\linewidth\linewidth\else\Gin@nat@width\fi}
\def\maxheight{\ifdim\Gin@nat@height>\textheight\textheight\else\Gin@nat@height\fi}
\makeatother
% Scale images if necessary, so that they will not overflow the page
% margins by default, and it is still possible to overwrite the defaults
% using explicit options in \includegraphics[width, height, ...]{}
\setkeys{Gin}{width=\maxwidth,height=\maxheight,keepaspectratio}
% Set default figure placement to htbp
\makeatletter
\def\fps@figure{htbp}
\makeatother
\setlength{\emergencystretch}{3em} % prevent overfull lines
\providecommand{\tightlist}{%
  \setlength{\itemsep}{0pt}\setlength{\parskip}{0pt}}
\setcounter{secnumdepth}{5}
% Workaround for mathspec clashes.
\makeatletter % undo the wrong changes made by mathspec
\let\RequirePackage\original@RequirePackage
\let\usepackage\RequirePackage
\makeatother

% Define options for mictrotype without touching default template.
\PassOptionsToPackage{expansion=false,final}{microtype}

% Activate colours to use in the book.
%\usepackage[usenames,dvipsnames,svgnames,table]{xcolor}
\PassOptionsToPackage{usenames,dvipsnames,svgnames,table}{xcolor}

% COLOR DEFINITIONS ------------------------------------------------------------
\definecolor{smartgray}{gray}{0.55}
\definecolor{smartred}{HTML}{800000}
\definecolor{smartplum}{HTML}{47294F}
\definecolor{smartblue}{HTML}{193A6B}
\definecolor{smartorange}{HTML}{AA4C00}
\definecolor{smartgreen}{HTML}{3C7705}
\definecolor{smartbutter}{HTML}{A08300}

\definecolor{tangored}{HTML}{CC0000}
\definecolor{tangoplum}{HTML}{75507B}
\definecolor{tangoblue}{HTML}{3465A4}
\definecolor{tangoorange}{HTML}{F57900}
\definecolor{tangogreen}{HTML}{73D216}
\definecolor{tangobutter}{HTML}{EDD400}

\definecolor{shadecolor}{gray}{0.95}

\colorlet{heading}{smartgray}
\colorlet{ref}{smartblue}
\colorlet{cite}{smartblue}
\colorlet{url}{smartgreen}
\colorlet{caption}{black}
\colorlet{margincaption}{caption}

% MARGIN -----------------------------------------------------------------------

\usepackage{ragged2e} % for marginnote

% test - make figures float on clear pages?
\makeatletter
\@fpsep\textheight
\makeatother

% FONT OPTIONS -----------------------------------------------------------------
\usepackage{fontspec}
% \usepackage{MinionPro}
% \usepackage{mathpazo}
% \setsansfont{Myriad Pro}
% \setsansfont{CMU Sans Serif}
% \setmonofont{Consolas}
\setmonofont{PT Mono}
\setmathfont{Minion Pro}


% Define a new 'widespaced' font that is used in section headings
% or chapter titles. For backwards compatibility reasons this uses
% the old calling convention \newfontfamily\fontname[<options>{font}
% instead of the newer \newfontfamily\fontname{font}[<options>]
\newfontfamily\widespaced[%
	LetterSpace=15,%
	WordSpace={1.25, 1.25, 1.25}%
]{Minion Pro}%


\DeclareRobustCommand{\spacedallcaps}[1]{%
	{\widespaced\oldstylenums{\MakeTextUppercase{#1}}}%
}%
\DeclareRobustCommand{\spacedlowsmallcaps}[1]{%
	{\widespaced\oldstylenums{\textsc{\MakeTextLowercase{#1}}}}%
}%

% CHAPTER STYLES ---------------------------------------------------------------

% \chapterstyle{ell}
% \renewcommand*{\chapnumfont}{\normalfont\huge\bfseries}
% \renewcommand*{\chaptitlefont}{\normalfont\Huge}
% \chapterstyle{brotherton}

% FIGURES ----------------------------------------------------------------------
% make sure floats do not appear before references
\usepackage{flafter}

\captionstyle{\small}
\captionnamefont{\slshape\color{caption}}
\captiondelim{\null\newline}
\captionnamefont{\small\bfseries\color{smartred}}
\captiontitlefont{\small\linespread{.84}\selectfont}
% \captionstyle[\raggedright]{}
\normalcaptionwidth
% \captiontitlefinal{.}
\setlength\belowcaptionskip{2.5ex}

\newsubfloat{figure}% Create subfloat in figure environment
\newsubfloat{table}% Create subfloat in table environment

% https://robjhyndman.com/hyndsight/latex-floats/
% better looking floats?

\setcounter{topnumber}{2} % max no of floats at top of page
\setcounter{bottomnumber}{2} % max no of floats at bottom of page
\setcounter{totalnumber}{2} % maximum number of floats on a page
\renewcommand{\topfraction}{0.85}
\renewcommand{\bottomfraction}{0.85}
\renewcommand{\textfraction}{0.15}
\renewcommand{\floatpagefraction}{0.5}

% COLOPHON ---------------------------------------------------------------------
\makeatletter
\newcommand{\acknowledgementname}{Acknowledgements}
\newenvironment{acknowledgement}{%
  \setup@bstract
  \if@bsrunin\else
    \ifnumber@bs \num@bs \else
      \begin{\absnamepos}\abstractnamefont\acknowledgementname\end\absnamepos%
      \vspace{\abstitleskip}%
    \fi
  \fi
  \put@bsintoc%
  \begin{@bstr@ctlist}\if@bsrunin\@bsrunintitle\fi\abstracttextfont}%
  {\par\end{@bstr@ctlist}}

\newcommand{\smartcopyright}[1]{\begingroup%
	\pagebreak%
	\thispagestyle{empty}%
	\footnotesize%
	\par\vspace*{\fill}%
	\hspace{-1.1cm}%
	\parbox{10cm}{%
		\noindent%
		\@title \\[.2cm]
		Copyright \textcopyright\ \the\year, \@author\\[0.5cm]
		{#1}%
	}%
\endgroup}

\makeatother

% % MARGINNOTE -------------------------------------------------------------------
% % There are three commands for placing elements in the margin:
% %
% % \marginnote{<TEXT>}
% % Just like \marginpar but with correct formating
% %
% % \marginfig{<SHORT CAPTION>}{<FIGURE CONTENT>}{<LONG CAPTION>}{\label{<LABEL>}}
% % \margintbl{<SHORT CAPTION>}{<TABLE CONTENT>}{<LONG CAPTION>}{\label{<LABEL>}}
% % Places a figure/table in the margin
% % - <SHORT CAPTION> is the caption as it occurs in the list of figures/tables
% % - <FIGURE CONTENT>/<TABLE CONTENT> is the actual figure/table content
% % - <LONG CAPTION> is the caption below the figure
% % - <LABEL> is the label -- we decided to let you explicitly type "\label" here
% % as this allows LaTeX editors to index the label and provide auto completion.
% 
% % Define the "marginnote" command. For compatibility reasons
% % we do not override marginpar.
% \strictpagecheck
% \newcommand{\marginnote}[1]{\hspace*{0pt}\marginpar{
% 	\microtypesetup{disable}%
% 	\slshape\footnotesize%
% 	\vspace*{-0.725em}%
% 	\checkoddpage%
% 	\ifoddpage%
% 		\RaggedRight%
% 	\else%
% 		\RaggedLeft%
% 	\fi%
% 	{#1}%
% 	\microtypesetup{enable}%
% }\ignorespaces}
% 
% \newcommand{\marginfloat}[8]{\marginnote{{#1}%
% \par\refstepcounter{#4}\textcolor{margincaption}{#5 #6}: {#2}#3}%
% \addcontentsline{#7}{#4}{\protect\numberline {#6}{\ignorespaces #8}}\ignorespaces}
% 
% \newcommand{\marginfig}[4]{\marginfloat{#2}{#3}{#4}{figure}{Figure}{\thefigure}{lof}{#1}}
\newcommand{\margintbl}[4]{\marginfloat{#2}{#3}{#4}{table}{Table}{\thetable}{lot}{#1}}


% PAGE STYLE (headers, footers) ------------------------------------------------
% Deactivate capitalization of headers.
\nouppercaseheads

% Define the new "berlin" style.
\makepagestyle{berlin}

% Display the page number in the footer.
\makeevenfoot{berlin}{\thepage}{}{}
\makeoddfoot{berlin}{}{}{\thepage}

% Override the chapter style with the unfinished "berlin" style
% to eliminate the header and to make sure a consistent footer
% style is used.
\copypagestyle{chapter}{berlin}

% lots of stuff taken out here...
% Figure this out later, eh?!


% FOOTNOTES --------------------------------------------------------------------
\footmarkstyle{\textsc{#1}}
\setlength{\footmarkwidth}{-0.5em}
\setlength{\footmarksep}{0.5em}
\setlength{\footparindent}{0em}
\setlength{\skip\footins}{2cm}

% TITLEPAGE --------------------------------------------------------------------
\makeatletter

\newcommand\institution[1]{\renewcommand\@institution{#1}}
\newcommand\@institution{}

\newcommand{\applytolist}[3][{,\\\vspace{0.1cm}}]{%
	\def\nextitem{\def\nextitem{#1}}%
	\@for \el:=#3\do{\nextitem{#2{\el}}}%
}

\newcommand{\smarttitle}{
{\begingroup
\thispagestyle{empty}
\makeatletter
\center
\vspace*{22 mm}
% {\Huge \textcolor{tangoblue}{\@title}\\}
{\Huge {\@title}\\}
\vfill
{\Large \@author\\}
\vspace{5 mm}
{\Large \@date}
\vfill
{\large \itshape {Submitted in partial fulfilment of the\\ requirements for the degree of}} \\
\vskip 0.2in
{\large Doctor of Philosophy}
\vskip 0.2in
{\large \itshape{\@institution}}
\vfill
\makeatother
\endgroup}}

\makeatletter


% OTHER PACKAGES AND SETTINGS --------------------------------------------------

% Context sensitive quotation facilities. Recommended by babel and should be 
% loaded before babel.
\usepackage{csquotes} 

% Sets the language used. Essential for proper hyphenation. Translates key words
% like 'figure' or 'table'. Note that 'english' is American English.
\usepackage[english]{babel} 

% Provides math enviroments, math symbols and more things useful for math.
\usepackage{latexsym,amsmath,amssymb,amsthm,amscd} 

% Enhances amsmath and provides further mathematical tools.
\usepackage{mathtools} 

% Provides inclusion of graphics and a proper interface.
\usepackage{graphicx} 

% Loading this exclusively for the \degree symbol!!
\usepackage{gensymb}

% Provides control over list environments. Supersedes 'enumerate', which gives 
% enumerate environment an optional argument which determines the style in which
% the counter is printed.
\usepackage{enumitem} 

% Hack to make floatrow work with memoir.
\let\newfloat\undefined 

% Handles alignment of floats (figures), with centering as default.
\usepackage{floatrow} 

% Ugly hack that sometimes helps with otherwise missing spaces.
\usepackage{xspace} 

% Automatically locates references on other pages. Load before hyperref.
\usepackage{varioref} 

% Clickable references.
\hypersetup{
  pdfencoding=auto,
  pdfborderstyle={/S/U/W 1},
  linkbordercolor=ref,
  colorlinks,
  linkcolor=ref,
  citecolor=cite,
  urlcolor=url}

% Anchors links to the beginning of their respective floats. 
% Load after hyperref.
\usepackage[all]{hypcap}

% Name references automatically. Load after hyperref.
% \usepackage[noabbrev,capitalize]{cleveref}

\usepackage[toc,acronym]{glossaries} % Glossary.


\usepackage{pgfplots}
\pgfplotsset{compat=1.8}
\pgfplotscreateplotcyclelist{smart}{%
  tangoplum,semithick,every mark/.append style={fill=tangoplum!80!black},mark=*\\%
  tangogreen,semithick,every mark/.append style={fill=tangogreen!80!black},mark=square*\\%
  tangoorange,semithick,every mark/.append style={fill=tangoorange!80!black},mark=otimes*\\%
  tangoblue,semithick,mark=star\\%
  tangobutter,semithick,every mark/.append style={fill=tangobutter!80!black},mark=diamond*\\%
  tangored,semithick,every mark/.append style={solid,fill=tangored!80!black},mark=*\\%
}
\pgfplotsset{cycle list name=smart}


% KABLE and KABLE-EXTRA
% \usepackage{array}
\usepackage{float}
\floatstyle{plaintop}
% \floatplacement{table}{t!} % not working
% \floatplacement{figure}{H}
\restylefloat{table}
\usepackage[tableposition=top, font=small, margin=10pt]{caption}
\usepackage{booktabs}
\usepackage{longtable}
\usepackage{array}
\usepackage{multirow}
\usepackage{wrapfig}
\usepackage{float}
\usepackage{colortbl}
\usepackage{pdflscape}
\usepackage{tabu}
\usepackage{threeparttable}
\usepackage{threeparttablex}
\usepackage[normalem]{ulem}
\usepackage{makecell}

\usepackage{bookmark} % 

% APPENDIX ---------------------------------------------------------------------

% load appendix package
\usepackage{appendix}

% PAGE LAYOUT ------------------------------------------------------------------
% \settypeblocksize{23.45cm}{11.85cm}{*} % original settings
\settypeblocksize{23.45cm}{13.05cm}{*} % expanded, since we are not using margin
% \setlength{\textheight}{23.45cm}
% \setlength{\textwidth}{11.85cm}
\setlength{\parskip}{1.1em} % add space between paragraphs
% \renewcommand{\baselinestretch}{1.2} % redefine line spacing
\setlrmargins{3.141592653589793cm}{*}{*}
\setulmargins{*}{*}{*}
\setmarginnotes{0.8cm}{3.5cm}{1cm}
\checkandfixthelayout

% let's not use hypenation in the text if justified:
\pretolerance=5000
\tolerance=9000
\emergencystretch=0pt
\righthyphenmin=4
\lefthyphenmin=4

% Ok, now for chapter style
% \chapterstyle{ntglike} % this is ok
% \chapterstyle{ger}
\renewcommand*{\chaptitlefont}{\LARGE \bfseries}
% \renewcommand*{\afterchaptertitle}%
    % {\newpage}

% \renewcommand{\chapterautorefname}{Chapter}
% \AtBeginDocument{\def\chapterautorefname{Chapter}}%

\setlength{\parindent}{0em} % don't indent paragraphs

% set numbering of subsections
\setsecnumdepth{subsubsection}

% set depth of table of contents
\setcounter{tocdepth}{2}

% prevent placing floats before a section (this should prevent figures from
% bunching at the end of the chapter.
% \usepackage[section]{placeins}

% \usepackage{tcolorbox} % for boxed text -- used in statement page.

% capitalise chapter cross reference
\renewcommand*{\chapterautorefname}{Chapter}

% inline code grey background hack
% using \codex activates highlighting of the code
% https://tex.stackexchange.com/a/140175
\newcommand\codex[2][]{
    \tikz[baseline=(s.base)]{
        \node(s)[
            rounded corners,
            fill=blue!5,        % background color
            % draw=gray,          % border of box
            % text=gray!50!black, % text color
            inner xsep =0pt,    % horizontal space between text and border
            inner ysep =0pt,    % vertical space between text and border
            text height=2ex,    % height of box
            text depth =1ex,    % depth of box
            #1                  % other options
        ]{\texttt{#2}};
    }
}


\usepackage{pdflscape}
\usepackage{afterpage}

\usepackage{siunitx}


\let\svthefootnote\thefootnote
\newcommand\blankfootnote[1]{%
  \let\thefootnote\relax\footnotetext{#1}%
  \let\thefootnote\svthefootnote%
}
\let\svfootnote\footnote
\renewcommand\footnote[2][?]{%
  \if\relax#1\relax%
    \blankfootnote{#2}%
  \else%
    \if?#1\svfootnote{#2}\else\svfootnote[#1]{#2}\fi%
  \fi
}
% Contains some macros which are commonly used in a thesis.
% Note that these are optional.
% Abbreviations
\newcommand*{\eg}{e.\,g.\@\xspace}
\newcommand*{\ie}{i.\,e.\@\xspace}
\newcommand*{\cf}{cf.\@\xspace}
\newcommand*{\etal}{et.\@ al.\@\xspace}
%\newcommand*{\etc}{etc.\@\xspace}
\makeatletter
\newcommand*{\etc}{%
  \@ifnextchar{.}%
    {etc}%
      {etc.\@\xspace}%
}
\makeatother

% Sets
\newcommand*{\R}{\mathbb{R}}
\newcommand*{\Q}{\mathbb{Q}}
\newcommand*{\Z}{\mathbb{Z}}
\newcommand*{\N}{\mathbb{N}}

% Matrix/vector operations
\newcommand*{\transpose}[1]{{#1^{\top}}}
\newcommand*{\inverse}[1]{{#1^{-1}}}
\newcommand*{\conjugate}[1]{{#1^{\ast}}}
\newcommand*{\pseudoinverse}[1]{{#1^{+}}}

% Stochastics
\newcommand*{\probability}[1]{{\mathrm{Pr}{#1}}}
\newcommand*{\expectation}[1]{{#1}}

% Optimization
\DeclareMathOperator*{\argmax}{arg\,max}
\DeclareMathOperator*{\argmin}{arg\,min}
\DeclareMathOperator{\lb}{lb}
\DeclareMathOperator{\sign}{sgn}

% Starts creation of a glossary.
\makeglossaries
\newglossaryentry{mjollnir}{
  name=Mjölnir,
  description={the hammer of Thor}}

\newglossaryentry{wolf}{
  name=wolf,
  plural=wolves,
  description={an awesome animal and a subspecies of \emph{Canis lupus}}}
  
\newglossaryentry{degc}% the label 
{% 
  name={\degree C},% the term 
  description={the set of all things},% a brief description 
} 


% Define all the metadata to be listed in \smarttitle and \smartcopyright
\title{Physiological response of adult and juvenile sea urchins (\textit{Heliocidaris erythrogramma}) to warming and acidification}
\author{Januar Harianto}
\institution{Discipline of Anatomy and Histology\\ The University of Sydney\\ 
  New South Wales, Australia}
\date{September, 2018}


% The integrated effects of projected climate change on cotton growth and physiology.

\author{Januar Harianto}
\date{}

\begin{document}

\frontmatter

\newgeometry{margin=1.8in}
\smarttitle
\restoregeometry

\SingleSpacing

\smartcopyright{This document is typeset in R Markdown and \LaTeX. The main typeface used is Minion Pro. Sans-serif text is typeset using \textsf{Myriad Pro} while source code is typeset using \texttt{PT Mono}. R code, plots and tables are mostly prepared in RStudio and Sketch on macOS.}

\OnehalfSpacing

\clearpage

\SingleSpacing

% \hrulefill

\section*{Author Attribution}

The work contained in the body of this thesis, except otherwise acknowledged, is the result of my own investigations. 

\textbf{Chapter \ref{chapter2}} is \href{https://doi.org/10/gdqdv4}{published} in Marine Biology. 
Januar Harianto (JH) designed and conducted the experiment with help from Sebastian Holmes (SH) and Maria Byrne (MB).
Hong Dao Nguyen (HDN) assisted with the collection of samples for protein analysis and SH helped with the experimental design of respirometry analyses. 
The manuscript was written by JH but all co-authors contributed intellectually to the drafts.

\textbf{Chapter \ref{chapter3}} is being prepared for submission. 
JH designed and conducted the experiment. 
The co-authors of this study include Sergio Torres-Gabarda (STG), Joshua Aldridge (JA), Richard Grainger (RG), MB, and JH. 
RG contributed to the collection of feeding rate data. 
All co-authors provided important feedback to manuscript drafts.

\textbf{Chapter \ref{chapter4}} is being prepared for submission.
The co-authors of this study are JH, STG, JA and MB. 
JH designed the experiment and conducted all the analyses. 
STG and JA contributed to the construction of the larval incubation system.
MB assisted in the collection of data from the larvae/juveniles.
All co-authors have reviewed drafts for the manuscript in preparation for submission.

\textbf{Chapter \ref{chapter5}} is published in Methods in Ecology and Evolution.
Nicholas Carey (NC) and JH co-wrote the manuscript. 
JH wrote the R code and code documentation. 
NC tested data and fixed errors.
NC and JH co-wrote the 15 vignettes for the package. 
MB provided vital contribution to the final manuscript.

\vfill
\hrulefill

% \noindent\fbox{%
%     \parbox{\textwidth}{%

% \begin{tcolorbox}
\section*{Thesis Statement of Originality}

This thesis has not been submitted for any other degree or purposes. I certify that the intellectual content of this thesis is the product of my own work and that all the assistance received in preparing this thesis and sources have been acknowledged.

\vspace{12mm}
\textbf{Januar Harianto}, 28 September, 2018

\vspace{3.14mm}

As supervisor for the candidature upon which this thesis is based, I can confirm that the authorship attribution statements above are correct.

\vspace{12mm}
\textbf{Maria Byrne}, 28 September, 2018

% \hrulefill

% \end{tcolorbox}


%   }%
% }
\cleardoublepage

\SingleSpacing

\chapter{Abstract}

\OnehalfSpacing

Approximately 25 \% to 40 \% of anthropogenic carbon dioxide (CO\textsubscript{2}) emissions have been absorbed by the oceans, influencing the physical and chemical properties of seawater and contributing to ocean warming (OW) and acidification (OA).
% These environmental stressors interact to affect biochemical reaction rates and physiological functioning in marine biota, thus expediting or mitigating biological impacts.
Marine organisms will have to adjust their physiology and behaviour to cope with these changes.
% However, whether they are able to maintain their fitness in the face of future climate change is dependent on the severity and time scale of the stressors involved.
This thesis investigated the impact of long-term exposure to OW and OA, at levels predicted for up to the year 2100 on the physiological responses of the sea urchin, \textit{Heliocidaris erythrogramma}. 
This species is an ecologically important benthic ecosystem engineer in south-east Australia with the ability to reshape community assemblages through its herbivory.
Experimental studies on adult \textit{H. erythrogramma} were conducted using experimental designs which involved a gradual introduction to future OW and OA conditions,  long-term incubation periods to accommodate acclimatory responses, and temperature adjustments over time to incorporate seasonal change. The influence of parental acclimation effects on offspring responses to OW and OA were also investigated to determine the potential for transgenerational effects.

Temperature is the most important factor in determining the performance and fitness of marine ectotherms. 
In the southeast Australian OW hotspot, up to +6 \degree C of warming is projected for the region by the year 2100. 
The consequences of exposure to warming were investigated in \textit{H. erythrogramma} in experiments that employed gradual adjustment over 6 weeks to four temperature levels (ambient 20 \degree C, 22 \degree C, 24 \degree C, 26 \degree C). This was followed by an acclimation period of 3 months at these temperatures in constant conditions.
Metabolic rate, $Q_{10}$, heat-shock protein (HSP70) expression, gonad index, gonad histology and survival were measured. 
There were no significant effects of elevated temperature on metabolic rate, HSP70 protein expression and survival at 22 \degree C, indicating that \textit{H. erythrogramma} was able to tolerate a +2 \degree C increase in temperature.
However, metabolic depression was evident at the highest warming level of +6 \degree C (26 \degree C) which caused decreased survival and elevated HSP70 protein expression.
\textit{Heliocidaris erythrogramma} responded to the stress of increased temperature by altering its physiology, but long-term exposure to +4 \degree C and +6 \degree C levels of warming incurred unsustainable metabolic costs and led to trade-offs between metabolism, HSP70 expression and survival. 

Decreased pH through uptake of CO\textsubscript{2} affects a wide range of physiological functions in sea urchins due to alteration of the carbonate chemistry of seawater which affects pH balance, mineral saturation for skeletogenesis and physiological hypercapnia (\textit{p}CO\textsubscript{2}).
The thermal and metabolic responses of \textit{H. erythrogramma} are likely to be influenced by OA, with the combined impact of OW and OA being difficult to predict.
To investigate the effects of both stressors on the physiology of \textit{H. erythrogramma}, OA was incorporated with OW in a multistressor experimental design.
\textit{Heliocidaris erythrogramma} collected in winter were gradually adjusted over 7 weeks to three temperature (ambient, +2 \degree C and +3 \degree C) and two pH (ambient 8.0, low 7.6) treatments.
Temperature was adjusted weekly to reflect the seasonal temperature change that occurs over time from winter to summer (17 \degree C to 22 \degree C). 
The elevated temperature treatments were adjusted to offset the temperature profile.
The sea urchins were incubated in those conditions for 22 weeks.

Metabolic rate was quantified at 4 and 12 weeks of acclimation and feeding and ammonia excretion rates were measured at 12 weeks. 
Survivorship was recorded daily for up to 22 weeks.
In single stressor treatments metabolic rate increased with either OW or OA.
In contrast, the two stressors interacted when combined causing a decrease in metabolism as temperature increases at low pH.
There was no evidence of interaction between the two stressors for feeding rate, ammonia excretion rate, or survival. 
Thus, while pH influenced metabolic rate causing metabolic depression at +3 \degree C, there was no evidence of pH affecting the other parameters measured. 
Survival was significantly reduced at +3 \degree C, regardless of pH, but not at +2 \degree C. 
% It appeared that the earlier arrival of warmer temperature at +3 \degree C was stressful, and the effect of this 1 \degree C difference was evident due to the inclusion of seasonal variation in the experiment.
It appeared that the earlier arrival of warmer temperature in the +3 \degree C treatment was stressful, especially during a time of increasing gonad development. 
It appears that \textit{H. erythrogramma} may be living close to its lethal limit in Sydney Harbour since it already experiences the temperatures used in the study.
This study highlighted the importance of investigating multiple physiological traits in characterising organism responses to climate change.
Current populations are likely to be vulnerable to the combined stress of OW and OA in the near future.

While investigations of climate change effects at the adult stage provide insights into the resilience of established adult populations, it is also important to understand how their offspring may be affected.
% Long-term studies may provide crucial information on the impacts of climate change stressors on the resilience or current populations of organisms, but the potential for adaptation (or acclimation) over time may also require an understanding of how responses may be reshaped over generations.
These transgenerational carry-over effects are influenced by a combination of parental environmental history and developmental plasticity and can convey resilience to offspring between generations.
To investigate the influence of parental environment on offspring performance, adult \textit{H. erythrogramma} were acclimated to OW and OA conditions over the gonad developmental period (winter to summer) and then spawned after 3 months in the adult treatments.
The resulting embryos were incubated over 14 days to the juvenile stage at four temperatures (18 \degree C, 20 \degree C, 22 \degree C, 24 \degree C) and two pH\textsubscript{T} levels (8.0, 7.6).
Juvenile metabolic rate and test diameter (as a proxy for growth) were measured.
Parental exposure to warming had a positive effect on the metabolic rate responses of juveniles to temperature, whereby they respired more (up to 30 \%) when incubated in warmer treatments.
There were no effects of parental environment or juvenile thermal environment on growth.
Parental exposure to low pH had no effect on metabolic rate or growth. 
However, exposure to low pH during offspring development significantly affected the growth of the juveniles (up to 4 \% decrease).
% Further research is necessary to disentangle multi-generational trans-generational effects and how they will facilitate persistence in a rapidly changing ocean.

% All of the studies performed in this thesis used metabolic rate as a physiological indicator of performance. 
This thesis uses metabolic rate as a key indicator of physiological response to projected climate change conditions.
Metabolic rate measurements are well-established as a key parameter to assess the biological and ecological consequences of environmental stress. 
However, current methods of estimating metabolic rates from respirometry data are difficult and are performed manually, which introduces subjectivity to the estimates.
To address problems associated with the lack of standard methods to analyse increasingly complex respirometry data, a new package for the R statistical computing language, called \texttt{respR}, was developed.
This software package provides a range of functions for analysing respirometry data.
Importantly, this package simplifies handling and analysis of respirometry data and can be used to reliably and rapidly generate reproducible analyses of metabolic rates.
The \texttt{respR} package provides new methods to accurately detect maximum and minimum rates, linear sections of previously intractable data, and critical oxygen tension, based novel techniques of rolling regression and machine learning.
The package has been published as an open-source project to address data reproducibility and transparency when processing respirometry data.


% This thesis provides evidence that the physiological and biological responses of animals to climate change stressors are influenced by the interactions between stressors, seasonal changes and trans-generational effects.
For \textit{H. erythrogramma}, an ecologically important species on Australia's temperate reefs, there were clear long-term physiological consequences of acclimation to OW and OA on physiological condition and fitness. 
This species is vulnerable to a +3 \degree C increase in temperature but appears resilient to decreased pH.
The response of \textit{H. erythrogramma} to the combined effect of both stressors appeared to be dominated by changes in temperature.
The outcomes of this research have significant implications as to the resilience of \textit{H. erythrogramma} to climate warming, especially where it resides in the global warming hotspot of temperate Australia.
However, there was evidence of carry-over effects on juvenile physiology shaped by parental environment, which suggests a capacity for transgenerational acclimation and resilience. 
Further work is required to determine the multi-generational response of this species to climate change.

% Lastly, data reproducibility was addressed with the introduction of an R package for respirometry which will allow physiologists to process respirometry data transparently.


% Overall this thesis shows that the sea urchin \textit{H. erythrogramma} was vulnerable to near future OW and OA, but evidence of carry-over effects on juvenile metabolism, shaped by parental environment, suggests that further work is required to determine the resilience of this species to climate change. 
% In addition, a novel computing package was developed which will hopefully allow future users to handle and process extensive respirometry datasets in a manner that is fully transparent and reproducible.

\clearpage


\SingleSpacing

\chapter*{Acknowledgements}

\OnehalfSpacing

This thesis, and the work contained within, would not have been possible without the wonderful support of an incredibly talented group of scholars and experts in their fields. I survived the process! All thanks to:

\begin{itemize}
  \item My supervisor, Dr. Maria Byrne. I am both honoured and grateful to be your student and collaborator. You are the most dedicated mentor that I have ever met. Thank you for all the opportunities that you have facilitated throughout these years. I am proud to have you as my supervisor, I also hope that I have, at the very least, met your expectations as a PhD student.
  \item Dr. Sebastian Holmes. Thanks, Seb, for believing in my work and encouraging me to push on. I am grateful for all the opportunities to work with you. Your unwavering commitment to get me through my final undergraduate year was the reason that I have not given up on academia!
  \item Dr. Hong Dao Nguyen. Without a doubt the most cheerful person I have ever met. Thanks for helping me get through the first year! And for being there for both friendship and advice. 
  \item Dr. Christina McGraw.  You made a programmer out of me, and I am still amazed that I can do so much with just a syringe pump, a spectrometer, and bunch of tubes. Thank you for your patience and kindness!
  % who is one of the kindest people I have ever met. She spent an entire week coaching me on instrument programming in LabVIEW, which got me really interested in making and building my own scientific devices later on.
  \item Dr. Jon Havenhand. You probably don't know it, but your introduction to R was one of the main reasons why I started using the programming language, and your encouraging words on my Excel template on analysing respirometry led to the decision to develop the R package that is now a chapter in this book. Thank you!
  % introduced me to the R programming language in Friday Harbour, and I would never have imagined myself writing functions and creating 
  \item Dr. Nicholas Carey. You made a better physiologist out of me, and chatting with you is always great fun. Thanks for the great advice on respirometry! I can call myself a physiologist because of you.
  \item Dr. Steven Doo. One day, Steve decided to force himself to learn R by re-analysing all his data using the programming language. \textit{What a great idea}, I thought! So I wrote my entire thesis in R. Didn't expect that, but thanks Steve.
  \item Sergio Torres-Gabarda. I couldn't have asked for a better lab mate and friend. You've really made the PhD experience that much easier by being there, sometimes magically, every time when I needed it. It's like having extra hands, because you're that good at reading my mind sometimes...did that mean I had 4 hands?
  \item Joshua Aldridge. Josh is probably the real reason why I was able to set up my experiments. He taught me everything I needed to know on how to build and maintain aquarium setups at SIMS. Basically, Josh turned me into a decent plumber. Josh is amazing. You're a great colleague and friend, Josh.
  \item Demian Koop and Paula Cisternas. Demian, thanks for the wonderful PhD advice throughout all these years. Paula, you have been a great lunch buddy -- thanks for all the advice over lunch and ramen, I have become a better person because of you.
  \item Roberta Johnson, Richard Grainger and Dione Deaker. These Honours students are now doing their PhDs, but back then they were my volunteers. They're the reason some people think that I can clone myself and do so many things in a day. Thanks guys! Thanks Roberta for working with me on all the crazy new ideas. We probably won't forget about those blood-gas analysers for a while. Rich, thanks for all the hard work with the feeding experiments. I would have not been able to do them without you. Dione, thank you for looking at the sea urchin spines! And for helping me throughout the end of the PhD with the random bits and pieces.
  \item People at the Sydney Institute of Marine Science, thank you for the opportunity to work at the aquarium facility and for accomodating my increasingly complex experimental setups. Amanda, Ulysse, Mark and Peter, SIMS is an amazing melting pot for collaboration thanks to you guys! 
  \item The crazy cool guys at CMEG. Augustine Porter. Marcelo Reis. Roberta Pereira. Elliot. Mitch. You made the experience in the building fun. Gus and Cello, we made it!
\end{itemize}

Special thanks to Kevin Keay and Frank Lovicu, who helped with various PhD matters and funding. Thanks to the Professor NGW \& Mrs Ann Macintosh Memorial Scholarship for financial support.

Last but not least, I would like to thank my wife, Cynthia, and my parents, Harianto and Lian Hiong for supporting me relentlessly throughout these years. I dedicate this thesis to my family.


\clearpage


\cleardoublepage

\SingleSpacing

\chapter{Preface}

\OnehalfSpacing

The data chapters presented in this thesis are written in a format for publication in peer-reviewed, scientific journals. 
Because they have been prepared as stand-alone manuscripts, repetition of important information was unavoidable and some chapter introductions may appear redundant to the General Introduction (\autoref{chapter1}).
In addition, all references have been re-formatted and compiled at the end of the thesis for easier cross-referencing.

This thesis has been saved as a PDF file and is submitted digitally. I have enabled hyperlinks in the document for the table of contents, citations, cross-references, geographic locations and the bibliography. 
These links are obvious -- they are coloured. Links in {\color{ref} blue} are cross-references within the document; links in {\color{url} green} point to websites, and will open in a web browser.
Clicking on any section heading in the table of contents will bring the reader to the relevant page, likewise, clicking on a citation will bring the reader directly to the relevant reference.
Where possible, every reference has a DOI link that can be clicked to access the online version of that reference through a web browser.
Clicking on a figure or table cross-reference links directly to the page containing that figure/table, and clicking on a GPS coordinate will automatically launch a web browser pinning that location on Google Maps.

The table of contents feature is fully supported. In Adobe Acrobat, displaying the Bookmarks toolbar will allow the reader to jump to specific sections or subsections of this thesis. On macOS, the default Preview app is also able to list the table of contents in the sidebar for easy navigation. If the reader is using a PDF reader with ``back/forward'' functionality, activating the ``back'' function will always bring the reader back to the same page before a link was clicked.
I hope that these small, but useful, additions to the PDF will help any reader(s) to better navigate this thesis.
There are probably dozens of you.

\mainmatter

% \raggedyright
\setlength{\parindent}{0em} % don't indent paragraphs

\hypersetup{linkcolor=black} % quick ugly hack for black toc

{
\setcounter{tocdepth}{1}
\tableofcontents
}
\hypertarget{biblio-style-apalike}{%
\chapter{biblio-style: apalike}\label{biblio-style-apalike}}

Placeholder

\hypertarget{the-consequences-of-climate-change}{%
\section{The consequences of climate change}\label{the-consequences-of-climate-change}}

\hypertarget{predicting-resilience-to-environmental-stress-through-investigation-of-physiology}{%
\section{Predicting resilience to environmental stress through investigation of physiology}\label{predicting-resilience-to-environmental-stress-through-investigation-of-physiology}}

\hypertarget{metabolic-rate-the-most-fundamental-bioindicator}{%
\section{Metabolic rate -- the most fundamental bioindicator}\label{metabolic-rate-the-most-fundamental-bioindicator}}

\hypertarget{survivorship}{%
\section{Survivorship}\label{survivorship}}

\hypertarget{key-anthropogenic-drivers-of-marine-physiological-function}{%
\section{Key anthropogenic drivers of marine physiological function}\label{key-anthropogenic-drivers-of-marine-physiological-function}}

\hypertarget{increased-temperature}{%
\subsection{Increased temperature}\label{increased-temperature}}

\hypertarget{increased-carbon-dioxide-concentration-pco2}{%
\subsection{\texorpdfstring{Increased carbon dioxide concentration, \emph{p}CO\textsubscript{2}}{Increased carbon dioxide concentration, pCO2}}\label{increased-carbon-dioxide-concentration-pco2}}

\hypertarget{on-the-complexities-of-climate-change-manipulation-experiments}{%
\section{On the complexities of climate change manipulation experiments}\label{on-the-complexities-of-climate-change-manipulation-experiments}}

\hypertarget{study-animal}{%
\section{Study animal}\label{study-animal}}

\hypertarget{effects-of-warming-and-acidification-on-h.-erythrogramma}{%
\subsection{\texorpdfstring{Effects of warming and acidification on H. \emph{erythrogramma}}{Effects of warming and acidification on H. erythrogramma}}\label{effects-of-warming-and-acidification-on-h.-erythrogramma}}

\hypertarget{thesis-aims-and-outline}{%
\section{Thesis aims and outline}\label{thesis-aims-and-outline}}

\hypertarget{chapter2}{%
\chapter{\texorpdfstring{The effect of warming on mortality, metabolic rate, heat-shock protein response and gonad growth in thermally acclimated sea urchins (\emph{Heliocidaris erythrogramma})}{The effect of warming on mortality, metabolic rate, heat-shock protein response and gonad growth in thermally acclimated sea urchins (Heliocidaris erythrogramma)}}\label{chapter2}}

Placeholder

\hypertarget{abstract}{%
\section*{Abstract}\label{abstract}}
\addcontentsline{toc}{section}{Abstract}

\hypertarget{introduction}{%
\section{Introduction}\label{introduction}}

\hypertarget{materials-and-methods}{%
\section{Materials and methods}\label{materials-and-methods}}

\hypertarget{collection}{%
\subsection{Collection}\label{collection}}

\hypertarget{habitat-temperature-profile}{%
\subsection{Habitat temperature profile}\label{habitat-temperature-profile}}

\hypertarget{acclimation-setup-and-maintenance}{%
\subsection{Acclimation setup and maintenance}\label{acclimation-setup-and-maintenance}}

\hypertarget{survival}{%
\subsection{Survival}\label{survival}}

\hypertarget{metabolic-rate}{%
\subsection{Metabolic rate}\label{metabolic-rate}}

\hypertarget{gonad-index-and-hsp70-expression}{%
\subsection{Gonad index and HSP70 expression}\label{gonad-index-and-hsp70-expression}}

\hypertarget{gonad-histology}{%
\subsection{Gonad Histology}\label{gonad-histology}}

\hypertarget{data-analyses}{%
\subsection{Data analyses}\label{data-analyses}}

\hypertarget{results}{%
\section{Results}\label{results}}

\hypertarget{habitat-temperature-profile-1}{%
\subsection{Habitat temperature profile}\label{habitat-temperature-profile-1}}

\hypertarget{survival-1}{%
\subsection{Survival}\label{survival-1}}

\hypertarget{metabolic-rate-1}{%
\subsection{Metabolic Rate}\label{metabolic-rate-1}}

\hypertarget{hsp70-expression}{%
\subsection{HSP70 expression}\label{hsp70-expression}}

\hypertarget{gonad-index}{%
\subsection{Gonad index}\label{gonad-index}}

\hypertarget{discussion}{%
\section{Discussion}\label{discussion}}

\hypertarget{chapter3}{%
\chapter{\texorpdfstring{Long-term study of the multistressor impacts of ocean warming and acidification on the physiology of the sea urchin, \emph{Heliocidaris erythrogramma}}{Long-term study of the multistressor impacts of ocean warming and acidification on the physiology of the sea urchin, Heliocidaris erythrogramma}}\label{chapter3}}

Placeholder

\hypertarget{abstract-1}{%
\section*{Abstract}\label{abstract-1}}
\addcontentsline{toc}{section}{Abstract}

\hypertarget{introduction-1}{%
\section{Introduction}\label{introduction-1}}

\hypertarget{methods}{%
\section{Methods}\label{methods}}

\hypertarget{collection-1}{%
\subsection{Collection}\label{collection-1}}

\hypertarget{experimental-setup}{%
\subsection{Experimental setup}\label{experimental-setup}}

\hypertarget{gradual-adjustment-to-treatment-conditions-and-acclimation}{%
\subsection{Gradual adjustment to treatment conditions and acclimation}\label{gradual-adjustment-to-treatment-conditions-and-acclimation}}

\hypertarget{monitoring-and-survival}{%
\subsection{Monitoring and survival}\label{monitoring-and-survival}}

\hypertarget{maintenance}{%
\subsection{Maintenance}\label{maintenance}}

\hypertarget{metabolic-rate-2}{%
\subsection{Metabolic rate}\label{metabolic-rate-2}}

\hypertarget{feeding-rate-and-assimilation-efficiency}{%
\subsection{Feeding rate and assimilation efficiency}\label{feeding-rate-and-assimilation-efficiency}}

\hypertarget{ammonia-excretion}{%
\subsection{Ammonia excretion}\label{ammonia-excretion}}

\hypertarget{scope-for-growth-sfg-parameters}{%
\subsection{Scope for growth (SfG) parameters}\label{scope-for-growth-sfg-parameters}}

\hypertarget{statistical-analyses}{%
\subsection{Statistical analyses}\label{statistical-analyses}}

\hypertarget{results-1}{%
\section{Results}\label{results-1}}

\hypertarget{survival-2}{%
\subsection{Survival}\label{survival-2}}

\hypertarget{metabolic-rate-3}{%
\subsection{Metabolic rate}\label{metabolic-rate-3}}

\hypertarget{feeding-rate}{%
\subsection{Feeding rate}\label{feeding-rate}}

\hypertarget{ammonia-excretion-rate}{%
\subsection{Ammonia excretion rate}\label{ammonia-excretion-rate}}

\hypertarget{scope-for-growth}{%
\subsection{Scope for growth}\label{scope-for-growth}}

\hypertarget{spawning-activity}{%
\subsection{Spawning activity}\label{spawning-activity}}

\hypertarget{discussion-1}{%
\section{Discussion}\label{discussion-1}}

\hypertarget{effect-on-survival}{%
\subsection{Effect on survival}\label{effect-on-survival}}

\hypertarget{effect-on-metabolism}{%
\subsection{Effect on metabolism}\label{effect-on-metabolism}}

\hypertarget{long-term-multistressor-impacts}{%
\subsection{Long-term multistressor impacts}\label{long-term-multistressor-impacts}}

\hypertarget{conclusions}{%
\subsection{Conclusions}\label{conclusions}}

\hypertarget{chapter4}{%
\chapter{\texorpdfstring{Offspring metabolic performance to ocean warming and acidification are shaped by parental environment in the sea urchin, \emph{Heliocidaris erythrogramma}}{Offspring metabolic performance to ocean warming and acidification are shaped by parental environment in the sea urchin, Heliocidaris erythrogramma}}\label{chapter4}}

Placeholder

\hypertarget{abstract-2}{%
\section*{Abstract}\label{abstract-2}}
\addcontentsline{toc}{section}{Abstract}

\hypertarget{introduction-2}{%
\section{Introduction}\label{introduction-2}}

\hypertarget{methods-1}{%
\section{Methods}\label{methods-1}}

\hypertarget{collection-of-adults}{%
\subsection{Collection of adults}\label{collection-of-adults}}

\hypertarget{experimental-design}{%
\subsection{Experimental design}\label{experimental-design}}

\hypertarget{adult-treatment-setup}{%
\subsection{Adult treatment setup}\label{adult-treatment-setup}}

\hypertarget{larval-kreisel-setup}{%
\subsection{Larval kreisel setup}\label{larval-kreisel-setup}}

\hypertarget{spawning-and-incubation}{%
\subsection{Spawning and incubation}\label{spawning-and-incubation}}

\hypertarget{metabolic-rate-and-growth}{%
\subsection{Metabolic rate and growth}\label{metabolic-rate-and-growth}}

\hypertarget{data-analyses-1}{%
\subsection{Data analyses}\label{data-analyses-1}}

\hypertarget{results-2}{%
\section{Results}\label{results-2}}

\hypertarget{metabolic-rate-4}{%
\subsection{Metabolic Rate}\label{metabolic-rate-4}}

\hypertarget{growth}{%
\subsection{Growth}\label{growth}}

\hypertarget{discussion-2}{%
\section{Discussion}\label{discussion-2}}

\hypertarget{chapter5}{%
\chapter{respR -- an R package for the analysis of respirometry data}\label{chapter5}}

Placeholder

\hypertarget{abstract-3}{%
\section*{Abstract}\label{abstract-3}}
\addcontentsline{toc}{section}{Abstract}

\hypertarget{introduction-3}{%
\section{Introduction}\label{introduction-3}}

\hypertarget{package-overview}{%
\section{Package overview}\label{package-overview}}

\hypertarget{data-import-and-exploration}{%
\section{Data import and exploration}\label{data-import-and-exploration}}

\hypertarget{data-analysis}{%
\subsection{Data analysis}\label{data-analysis}}

\hypertarget{calculating-background-rate}{%
\subsubsection{Calculating background rate}\label{calculating-background-rate}}

\hypertarget{processing-respirometry-data}{%
\subsubsection{Processing respirometry data}\label{processing-respirometry-data}}

\hypertarget{estimating-rate-from-flow-through-respirometry}{%
\subsubsection{Estimating rate from flow-through respirometry}\label{estimating-rate-from-flow-through-respirometry}}

\hypertarget{estimating-critical-oxygen-tension-p_crit}{%
\subsubsection{\texorpdfstring{Estimating critical oxygen tension, P\(_{crit}\)}{Estimating critical oxygen tension, P\_\{crit\}}}\label{estimating-critical-oxygen-tension-p_crit}}

\hypertarget{correcting-converting-and-scaling-data}{%
\section{Correcting, converting and scaling data}\label{correcting-converting-and-scaling-data}}

\hypertarget{corrections-for-background-respiration}{%
\subsection{Corrections for background respiration}\label{corrections-for-background-respiration}}

\hypertarget{rate-conversions}{%
\subsection{Rate conversions}\label{rate-conversions}}

\hypertarget{assessing-performance-in-detecting-linear-data}{%
\section{Assessing performance in detecting linear data}\label{assessing-performance-in-detecting-linear-data}}

\hypertarget{generating-simulated-data-for-tests}{%
\subsection{Generating simulated data for tests}\label{generating-simulated-data-for-tests}}

\hypertarget{default-data}{%
\subsubsection{\texorpdfstring{\enquote{Default} data}{``Default'' data}}\label{default-data}}

\hypertarget{corrupted-data}{%
\subsubsection{\texorpdfstring{\enquote{Corrupted} data}{``Corrupted'' data}}\label{corrupted-data}}

\hypertarget{segmented-data}{%
\subsubsection{\texorpdfstring{\enquote{Segmented} data}{``Segmented'' data}}\label{segmented-data}}

\hypertarget{test-conditions}{%
\subsection{Test conditions}\label{test-conditions}}

\hypertarget{results-3}{%
\section{Results}\label{results-3}}

\hypertarget{default-data-1}{%
\subsection{\texorpdfstring{\enquote{Default} data}{``Default'' data}}\label{default-data-1}}

\hypertarget{corrupted-data-1}{%
\subsection{\texorpdfstring{\enquote{Corrupted} data}{``Corrupted'' data}}\label{corrupted-data-1}}

\hypertarget{segmented-data-1}{%
\subsection{\texorpdfstring{\enquote{Segmented} data}{``Segmented'' data}}\label{segmented-data-1}}

\hypertarget{concluding-remarks-and-future-improvements}{%
\section{Concluding remarks and future improvements}\label{concluding-remarks-and-future-improvements}}

\hypertarget{general-discussion}{%
\chapter{General discussion}\label{general-discussion}}

Placeholder

\hypertarget{metabolic-rate-is-a-key-physiological-response-variable}{%
\section{Metabolic rate is a key physiological response variable}\label{metabolic-rate-is-a-key-physiological-response-variable}}

\hypertarget{timing-and-magnitude-of-stress-events-can-influence-the-outcome-of-stressor-experiments}{%
\section{Timing and magnitude of stress events can influence the outcome of stressor experiments}\label{timing-and-magnitude-of-stress-events-can-influence-the-outcome-of-stressor-experiments}}

\hypertarget{offspring-responses-are-influenced-by-parental-environments}{%
\section{Offspring responses are influenced by parental environments}\label{offspring-responses-are-influenced-by-parental-environments}}

\hypertarget{modern-analytical-techniques-to-improve-the-analysis-respirometry-data}{%
\section{Modern analytical techniques to improve the analysis respirometry data}\label{modern-analytical-techniques-to-improve-the-analysis-respirometry-data}}

\hypertarget{concluding-remarks-and-future-directions}{%
\section{Concluding remarks and future directions}\label{concluding-remarks-and-future-directions}}

\end{document}
